\chapter*{Kata Pengantar}
\addcontentsline{toc}{chapter}{Kata Pengantar}
\markboth{Kata Pengantar}{Kata Pengantar}

Buku ajar ini disusun untuk mendukung pembelajaran mata kuliah
\textbf{Teknik Kompilasi} pada Program Studi Teknik Informatika,
Fakultas Teknologi Informasi, Universitas Bale Bandung. Mata kuliah ini
membahas bagaimana sebuah program sumber berbahasa tingkat tinggi
diterjemahkan menjadi representasi yang dapat diproses lebih lanjut
oleh mesin. Penekanan buku ini adalah \emph{front-end} kompilator:
analisis leksikal, analisis sintaksis, analisis semantik, serta
pembentukan dan validasi \emph{Abstract Syntax Tree} (AST).

Penyusunan buku mengikuti Rencana Pembelajaran Semester (RPS) berbasis
\emph{Outcome-Based Education} (OBE). Setiap bab dimulai dengan outcome
pembelajaran dan Sub-CPMK yang jelas, dilanjutkan dengan materi,
contoh, praktikum, latihan, asesmen, daftar cek kompetensi, rangkuman,
dan evaluasi. Dengan struktur ini, mahasiswa mengetahui sejak awal
kompetensi yang harus dicapai dan bukti belajar yang harus ditunjukkan.

Buku ini terdiri atas tujuh belas bab. Bab I dan Bab II memberikan
orientasi serta landasan teori. Bab III sampai Bab XVI merupakan empat
belas bab inti yang mengikuti empat belas topik RPS secara berurutan.
Bab XVII mengintegrasikan seluruh kompetensi melalui evaluasi,
refleksi, dan proyek akhir. Empat lampiran memuat rubrik penilaian,
contoh tugas dan portofolio, glosarium, serta sumber belajar tambahan.

Sepanjang semester, mahasiswa mengerjakan satu proyek berkelanjutan:
mini-kompilator untuk bahasa \textbf{Subset-C}. Proyek ini tumbuh dari
bab ke bab, mulai dari spesifikasi bahasa dan \emph{lexer}, dilanjutkan
dengan \emph{parser}, AST, tabel simbol, pemeriksaan tipe, hingga
validator dan diagnostik. Dengan cara ini, teori yang dipelajari
langsung berubah menjadi artefak perangkat lunak yang dapat diuji.

Buku ini memanfaatkan sumber-sumber terbuka yang telah tersedia,
antara lain karya Thain, Mogensen, Nystrom, Guntoro, dan catatan kuliah
Asep Juarna, serta sejumlah sumber daring yang tercantum pada daftar
pustaka. Penulis menyadari bahwa buku ajar ini masih dapat
disempurnakan. Kritik dan saran dari dosen, mahasiswa, dan pembaca
sangat diharapkan untuk perbaikan edisi berikutnya.

\begin{flushright}
Bandung, \BookYear\\
\BookAuthor
\end{flushright}
